\chapter{Control Flow}

\section{\texttt{if} / \texttt{else}}

\begin{lstlisting}[caption={Basic if/else structure}]
let grade = 7.5

if grade >= 7 {
  display "passed"
} else if grade >= 5 {
  display "make-up exam"
} else {
  display "failed"
}
\end{lstlisting}

\subsection*{\texttt{if} as an expression}

\lstinline{if} can be used as an expression — it returns the value of the
chosen branch:

\begin{lstlisting}
let status = if grade >= 7 { "passed" } else { "failed" }
display status
\end{lstlisting}

When used as an expression, both branches must have the same type (or both
\texttt{nil}).

\section{\texttt{match}}

\lstinline{match} compares a value against multiple patterns. The first
matching pattern is executed:

\begin{lstlisting}[caption={match with literals}]
let code = 2

match code {
  1 => display "beginner"
  2 => display "intermediate"
  3 => display "advanced"
  _ => display "unknown"
}
\end{lstlisting}

The \lstinline{_} pattern is the wildcard — it matches any value.

\subsection*{\texttt{match} with strings}

\begin{lstlisting}
let state = "TX"

match state {
  "TX" => display "Texas"
  "CA" => display "California"
  "NY" => display "New York"
  _    => display f"Unknown state: {state}"
}
\end{lstlisting}

\subsection*{\texttt{match} as an expression}

\begin{lstlisting}
let description = match code {
  1 => "beginner"
  2 => "intermediate"
  _ => "advanced"
}
\end{lstlisting}

\subsection*{Multiple values per arm}

\begin{lstlisting}
match code {
  1 | 2    => display "basic"
  3 | 4    => display "advanced"
  _        => display "other"
}
\end{lstlisting}

\section{Expression block}

A \lstinline|\{ \}| block is an expression that evaluates to the last
value produced inside it:

\begin{lstlisting}
let result = {
  let a = 10
  let b = 20
  a + b      // this is the block's value
}
display result    // 30
\end{lstlisting}
